
\subsection{忍耐}
\subsubsection{路加福音8:15/21:14-19}
\begin{itemize}[noitemsep]
\item 8:15 那落在好土里的，就是人听了道，持守在诚实善良的心里，并且\textcolor{red}{忍耐}着结实。
\item 21:14-19 所以，你们当立定心意，不要预先思想怎样分诉；因为我必赐你们口才、智慧，是你们一切敌人所敌不住、驳不倒的。连你们的父母、弟兄、亲族、朋友也要把你们交官，你们也有被他们害死的。你们要为我的名被众人恨恶，然而，你们连一根头发也必不损坏。你们常存\textcolor{red}{忍耐}，就必保全灵魂(必得生命)。
\end{itemize}
\subsubsection{罗马书2:7/5:3-4/8:25/15:4-5} 
\begin{itemize}[noitemsep]
\item 2:7 他必照各人的行为报应各人。凡\textcolor{red}{恒心(耐心)}行善寻求荣耀、尊贵和不能朽坏之福的,就以永生报应他们
\item 5:3-4 我们既因信称义，就借着我们的主耶稣基督得与神相和。 2 我们又借着他，因信得进入现在所站的这恩典中，并且欢欢喜喜盼望神的荣耀。 3 不但如此，就是在患难中也是欢欢喜喜的。因为知道患难生\textcolor{red}{忍耐}， 4 \textcolor{red}{忍耐}生老练，老练生盼望， 5 盼望不至于羞耻，因为所赐给我们的圣灵将神的爱浇灌在我们心里。
\item 8:25 我们得救是在乎盼望。只是所见的盼望不是盼望。谁还盼望他所见的呢？25 但我们若盼望那所不见的，就必\textcolor{red}{忍耐}等候。
\item 15:4-5 从前所写的圣经都是为教训我们写的，叫我们因圣经所生的\textcolor{red}{忍耐}和安慰，可以得着盼望。 5 但愿赐\textcolor{red}{忍耐}、安慰的神叫你们彼此同心，效法基督耶稣， 6 一心一口荣耀神我们主耶稣基督的父！
\end{itemize}

\subsubsection{哥林多后书 1:6/6:4/12:12} 
\begin{itemize}[noitemsep]
\item 1:6 我们在一切患难中，他就安慰我们，叫我们能用神所赐的安慰去安慰那遭各样患难的人。 5 我们既多受基督的苦楚，就靠基督多得安慰。 6 我们受患难呢，是为叫你们得安慰、得拯救；我们得安慰呢，也是为叫你们得安慰。这安慰能叫你们\textcolor{red}{忍受}我们所受的那样苦楚。 7 我们为你们所存的盼望是确定的，因为知道你们既是同受苦楚，也必同得安慰。
\item 6:4 我们凡事都不叫人有妨碍，免得这职分被人毁谤；反倒在各样的事上表明自己是神的用人，就如在许多的\textcolor{red}{忍耐}，患难，穷乏，困苦，鞭打，监禁，扰乱，勤劳，警醒，不食，廉洁，知识，恒忍，恩慈，圣灵的感化，无伪的爱心，真实的道理，神的大能；仁义的兵器在左在右，荣耀羞辱，恶名美名；似乎是诱惑人的，却是诚实的；似乎不为人所知，却是人所共知的；似乎要死，却是活着的；似乎受责罚，却是不致丧命的；似乎忧愁，却是常常快乐的；似乎贫穷，却是叫许多人富足的；似乎一无所有，却是样样都有的。
\item 12:12 我成了愚妄人，是被你们强逼的。我本该被你们称许才是。我虽算不了什么，却没有一件事在那些最大的使徒以下。 12 我在你们中间用百般的\textcolor{red}{忍耐}，借着神迹、奇事、异能，显出使徒的凭据来。 13 除了我不累着你们这一件事，你们还有什么事不及别的教会呢？这不公之处，求你们饶恕我吧！
\end{itemize}

\subsubsection{帖前 1:3/帖后 1:4/3:5} 
\begin{itemize}[noitemsep]
\item 帖前 1:3 我们为你们众人常常感谢神，祷告的时候提到你们， 3 在神我们的父面前不住地记念你们因信心所做的工夫，因爱心所受的劳苦，因盼望我们主耶稣基督所存的\textcolor{red}{忍耐}。
\item 帖后 1:4 弟兄们，我们该为你们常常感谢神，这本是合宜的，因你们的信心格外增长，并且你们众人彼此相爱的心也都充足。 4 甚至我们在神的各教会里为你们夸口，都因你们在所受的一切逼迫患难中，仍旧存\textcolor{red}{忍耐}和信心。
\item 3:5 主是信实的，要坚固你们，保护你们脱离那恶者(脫離凶惡)。 4 我们靠主深信，你们现在是遵行我们所吩咐的，后来也必要遵行。 5 愿主引导你们的心，叫你们爱神，并学基督的\textcolor{red}{忍耐}！
\end{itemize}

\subsubsection{提前 6:11/提后 3:10/提多书 2:2} 
\begin{itemize}[noitemsep]
\item 提前 6:11 那些想要发财的人，就陷在迷惑，落在网罗和许多无知有害的私欲里，叫人沉在败坏和灭亡中。 10 贪财是万恶之根。有人贪恋钱财，就被引诱离了真道，用许多愁苦把自己刺透了。但你这属神的人要逃避这些事，追求公义、敬虔、信心、爱心、\textcolor{red}{忍耐}、温柔。 
\item 提后 3:10 从前雅尼和佯庇怎样敌挡摩西，这等人也怎样敌挡真道。他们的心地坏了，在真道上是可废弃的。 9 然而他们不能再这样敌挡，因为他们的愚昧必在众人面前显露出来，像那二人一样。 10 但你已经服从了我的教训、品行、志向、信心、宽容、爱心、\textcolor{red}{忍耐}， 11 以及我在安提阿、以哥念、路司得所遭遇的逼迫、苦难。我所忍受是何等的逼迫，但从这一切苦难中，主都把我救出来了。 12 不但如此，凡立志在基督耶稣里敬虔度日的，也都要受逼迫。 13 只是作恶的和迷惑人的必越久越恶，他欺哄人，也被人欺哄。 14 但你所学习的、所确信的，要存在心里，因为你知道是跟谁学的， 15 并且知道你是从小明白圣经，这圣经能使你因信基督耶稣有得救的智慧。
\item 提多 2:2 劝老年人要有节制，端庄，自守，在信心、爱心、\textcolor{red}{忍耐}上都要纯全无疵。
\end{itemize}

\subsubsection{希伯来书 10:36/12:1} 
\begin{itemize}[noitemsep]
\item 10:36 你们要追念往日蒙了光照以后，所忍受大争战的各样苦难。一面被毁谤，遭患难，成了戏景叫众人观看；一面陪伴那些受这样苦难的人。因为你们体恤了那些被捆锁的人，并且你们的家业被人抢去，也甘心忍受，知道自己有更美、长存的家业。 所以你们不可丢弃勇敢的心，存这样的心必得大赏赐。你们必须\textcolor{red}{忍耐},使你们行完了神的旨意,就可以得着所应许的。 37 “因为还有一点点时候，那要来的就来，并不迟延。

\item 12:1 我们既有这许多的见证人，如同云彩围着我们，就当放下各样的重担，脱去容易缠累我们的罪，存心\textcolor{red}{忍耐}，奔那摆在我们前头的路程， 2 仰望为我们信心创始成终的耶稣(仰望那将真道创始成终的耶稣。)。他因那摆在前面的喜乐，就轻看羞辱，忍受了十字架的苦难，便坐在神宝座的右边。 3 那忍受罪人这样顶撞的，你们要思想，免得疲倦灰心。 4 你们与罪恶相争，还没有抵挡到流血的地步。 5 你们又忘了那劝你们如同劝儿子的话说：“我儿，你不可轻看主的管教，被他责备的时候也不可灰心。 6 因为主所爱的，他必管教，又鞭打凡所收纳的儿子。” 7 你们所忍受的，是神管教你们，待你们如同待儿子。焉有儿子不被父亲管教的呢？ 8 管教原是众子所共受的，你们若不受管教，就是私子，不是儿子了。
\end{itemize}
\subsubsection{雅各书 1:3/5:10-11} 
\begin{itemize}[noitemsep]
\item 1:3 我的弟兄们，你们落在百般试炼中，都要以为大喜乐， 3 因为知道，你们的信心经过试验就生\textcolor{red}{忍耐}。 4 但\textcolor{red}{忍耐}也当成功，使你们成全、完备，毫无缺欠。

\item 5:10-11 弟兄们，你们不要彼此埋怨，免得受审判。看哪，审判的主站在门前了！ 10 弟兄们，你们要把那先前奉主名说话的众先知当做能受苦、能\textcolor{red}{忍耐}的榜样。 11 那先前\textcolor{red}{忍耐}的人，我们称他们是有福的。你们听见过约伯的\textcolor{red}{忍耐}，也知道主给他的结局，明显主是满心怜悯、大有慈悲。
\end{itemize}
\subsubsection{歌罗西书 1:11/彼得后书 1:6/啟示錄 3:10-11} 
\begin{itemize}[noitemsep]
\item 歌罗西书 1:11 我们自从听见的日子，也就为你们不住地祷告祈求，愿你们在一切属灵的智慧悟性上，满心知道神的旨意， 10 好叫你们行事为人对得起主，凡事蒙他喜悦，在一切善事上结果子，渐渐地多知道神； 11 照他荣耀的权能，得以在各样的力上加力，好叫你们凡事欢欢喜喜地\textcolor{red}{忍耐}宽容。
 
\item 彼得后书 1:6 你们要分外地殷勤：有了信心，又要加上德行；有了德行，又要加上知识； 6 有了知识，又要加上节制；有了节制，又要加上\textcolor{red}{忍耐}；有了\textcolor{red}{忍耐}，又要加上虔敬； 7 有了虔敬，又要加上爱弟兄的心；有了爱弟兄的心，又要加上爱众人的心。

\item 啟示錄 3:10-11 你既遵守我\textcolor{red}{忍耐}的道，我必在普天下人受試煉的時候，保守你免去你的試煉。 11我必快來！你要持守你所有的，免得人奪去你的冠冕。
\end{itemize}









\subsection{神的榮耀}
榮耀 Glory，希伯來字 Kabowd ＃3519，來自字根 kabad 本意為「重」的意思，引申為價值，尊貴，榮美，豐富，榮譽；以西結在書中，用虹的光彩來描寫神的榮耀。神的榮耀就是祂的本體，彰顯祂的聖潔，榮美，能力，與莊嚴，是神本相的表現。所以詩人說「諸天述說神的榮耀」（詩 19：1），就是藉著神所創造的宇宙萬物，神啟示給我們看到祂的本體。
榮耀的希臘字是 doxa ＃1391，原意是「信念」belief 與「看法」opinion 的意思，由七十士譯本將希伯來文的 Kabowd 譯成希臘文的 doxa，引申為讚美，尊崇，榮耀 ；神的榮耀就是神所有屬性的總和。
%\subsubsection{以西結書}


%\subsubsection{出埃及記}

%摩西在出埃及記中求神彰顯榮耀給他看，神回答「，又說「以西結書從神對選民長久的忍耐（不輕易發怒），回轉的勸誡（有憐憫有恩典），猶大與外邦所受的刑罰（萬不以有罪的為無罪），到餘民復興的預言，聖殿的重建（為千萬人存留慈愛，赦免罪孽），無不彰顯神的榮耀（屬性）。

%\subsubsection{羅 3：23}
%保羅說「世人都犯了罪，虧缺了神的榮耀」（羅 3：23），也就是說我們因著罪的緣故，與神的性情，神的本質分隔。然而神的心意是要人回到神的榮耀中，
%\subsubsection{猶 1：24}
%「叫你們無瑕無疵，歡歡喜喜站在他榮耀之前…」（猶 1：24），
%\subsubsection{啟 21：23}
%「那城內又不用日月光照，因有神的榮耀光照…」（啟 21：23），唯一的道路就是主耶穌基督。
\subsubsection{結/約1:18/2:19－21/13:31/14:9/17:1－6}
\begin{table}[H]
\centering
\begin{tabular}{p{2.cm}|p{15.50cm}}\hline\hline
\multicolumn{1}{c|}{经文}&\multicolumn{1}{c}{内容}\\\hline
出34:6－7&耶和華在他面前宣告說，耶和華，耶和華，是有憐憫，有恩典的神，不輕易發怒，並有豐盛的慈愛和誠實。  為千萬人存留慈愛，赦免罪孽，過犯，和罪惡。萬不以有罪的為無罪，必追討他的罪…」\\\hline
出3:18－19&我要顯我一切的恩慈，在你面前經過，宣告我的名…」\\\hline

結&以西結在蒙召時得以見到神的面，看到祂的榮耀\\\hline
結&猶大國因長期悖逆叛道，在聖城被毀前，神的榮耀從城中離開，停在城東的山上\\\hline
結&神又給先知看到更新的聖城，神的榮耀再次在萬民中彰顯，「人就知道我是耶和華。」\\\hline
約1:18&從來沒有人看見神，只有在父懷裡的獨生子將祂表明出來\\\hline
約2:19－21&耶穌回答說：「你們拆毀這殿，我三日內要再建立起來。」 20 猶太人便說：「這殿是四十六年才造成的，你三日內就再建立起來嗎？」 21 但耶穌這話是以他的身體為殿。\\\hline
約13:31&如今人子得了榮耀，神在人子身上也得了榮耀\\\hline
約14:9&看見了我，就是看見了父\\\hline
約17:1－6&父阿，時候到了。願你榮耀你的兒子，使兒子也榮耀你。  正如你曾賜給他權柄，管理凡有血氣的，叫他將永生賜給你所賜給他的人。  認識你獨一的真神，並且認識你所差來的耶穌基督，這就是永生。我在地上已經榮耀你，你所託付我的事，我已成全了。父阿，現在求你使我同你享榮耀，就是未有世界以先，我同你所有的榮耀。  你從世上賜給我的人，我已將你的名顯明與他們。\\\hline
羅3:23&世人都犯了罪，虧缺了神的榮耀\\\hline
彼前2:12&你們在外邦人中，應當品行端正，叫那些毀謗你們是作惡的，因看見你們的好行為，便在鑒察的日子，歸榮耀給神。\\\hline
猶1:24&叫你們無瑕無疵，歡歡喜喜站在他榮耀之前…\\\hline
啟21:23&那城內又不用日月光照，因有神的榮耀光照…\\\hline

\end{tabular}
\end{table}








\subsection{屍首必給空中的飛鳥和地上的野獸做食物}
\subsubsection{耶利米書 34：17-22}
所以耶和華如此說：你們沒有聽從我，各人向弟兄、鄰舍宣告自由，看哪，我向你們宣告一樣自由，就是使你們自由於刀劍、饑荒、瘟疫之下，並且使你們在天下萬國中拋來拋去。這是耶和華說的。 18 猶大的首領，耶路撒冷的首領，太監、祭司和國中的眾民曾將牛犢劈開，分成兩半，從其中經過，在我面前立約。後來又違背我的約，不遵行這約上的話。 19  20 我必將他們交在仇敵和尋索其命的人手中，他們的屍首必給空中的飛鳥和地上的野獸做食物。 21 並且我必將猶大王西底家和他的首領，交在他們仇敵和尋索其命的人與那暫離你們而去巴比倫王軍隊的手中。 22 耶和華說：我必吩咐他們回到這城，攻打這城，將城攻取，用火焚燒。我也要使猶大的城邑變為荒場，無人居住。













  

\subsection{天上的星，海边的沙}

創世記 15:5
於是領他走到外邊，說：「你向天觀看，數算眾星，能數得過來嗎？」又對他說：「你的後裔將要如此。」

創世記 22:17 論福，我必賜大福給你；論子孫，我必叫你的子孫多起來，如同天上的星，海邊的沙；你子孫必得著仇敵的城門；

創世記 26:4
我要加增你的後裔，像天上的星那樣多，又要將這些地都賜給你的後裔，並且地上萬國必因你的後裔得福。

創世記 28:14
你的後裔必像地上的塵沙那樣多，必向東西南北開展。地上萬族必因你和你的後裔得福。

創世記 32:12
你曾說：『我必定厚待你，使你的後裔如同海邊的沙，多得不可勝數。』」

出埃及記 32:13
求你記念你的僕人亞伯拉罕、以撒、以色列，你曾指著自己起誓說：『我必使你們的後裔像天上的星那樣多，並且我所應許的這全地，必給你們的後裔，他們要永遠承受為業。』」

申命記 1:10
耶和華你們的神使你們多起來，看哪，你們今日像天上的星那樣多！
申命記 10:22
你的列祖七十人下埃及，現在耶和華你的神使你如同天上的星那樣多。

撒母耳記下 17:11 依我之計，不如將以色列眾人從但直到別是巴，如同海邊的沙那樣多聚集到你這裡來，你也親自率領他們出戰。

列王紀上 4:20
猶大人和以色列人如同海邊的沙那樣多，都吃喝快樂。
列王紀上 4:29
神賜給所羅門極大的智慧、聰明和廣大的心，如同海沙不可測量

歷代志上 27:23
以色列人二十歲以內的，大衛沒有記其數目，因耶和華曾應許說必加增以色列人如天上的星那樣多。
歷代志下 1:9
耶和華神啊，現在求你成就向我父大衛所應許的話，因你立我做這民的王，他們如同地上塵沙那樣多。

尼希米記 9:23 你也使他們的子孫多如天上的星，帶他們到你所應許他們列祖進入得為業之地。

以賽亞書 48:19 你的後裔也必多如海沙，你腹中所生的也必多如沙粒，他的名在我面前必不剪除，也不滅絕。」

耶利米書 33:22
天上的萬象不能數算，海邊的塵沙也不能斗量，我必照樣使我僕人大衛的後裔和侍奉我的利未人多起來。」

何西阿書 1:10
「然而，以色列的人數必如海沙，不可量，不可數。從前在什麼地方對他們說『你們不是我的子民』，將來在那裡必對他們說『你們是永生神的兒子』。


羅馬書 9:27
以賽亞指著以色列人喊著說：「以色列人雖多如海沙，得救的不過是剩下的餘數

希伯來書 11:12
所以從一個彷彿已死的人就生出子孫，如同天上的星那樣眾多，海邊的沙那樣無數。

啟示錄 12:17
龍向婦人發怒，去與她其餘的兒女爭戰，這兒女就是那守神誡命、為耶穌作見證的。那時龍就站在海邊的沙上。







\subsection{鐵杖權柄制伏列國——推雅推喇}
\subsubsection{詩篇 2:6-12，110:2}
\begin{itemize} 
\itemsep-.25em
\item 你求我，我就將列國賜你為基業，將地極賜你為田產。
7 受膏者說：「我要傳聖旨。耶和華曾對我說：『你是我的兒子，我今日生你。
8 你求我，我就將列國賜你為基業，將地極賜你為田產。
9 你必用鐵杖打破他們，你必將他們如同窯匠的瓦器摔碎。』」
10 現在你們君王應當醒悟，你們世上的審判官該受管教！
11 當存畏懼侍奉耶和華，又當存戰兢而快樂。
12 當以嘴親子，恐怕他發怒，你們便在道中滅亡，因為他的怒氣快要發作。凡投靠他的，都是有福的！
\item 耶和華必使你從錫安伸出能力的杖來，你要在你仇敵中掌權。
\end{itemize} 

\subsubsection{以賽亞書 11:1-6}
1 從耶西的本必發一條，從他根生的枝子必結果實。 2 耶和華的靈必住在他身上，就是使他有智慧和聰明的靈，謀略和能力的靈，知識和敬畏耶和華的靈。 3 他必以敬畏耶和華為樂，行審判不憑眼見，斷是非也不憑耳聞； 4 卻要以公義審判貧窮人，以正直判斷世上的謙卑人，以口中的杖擊打世界，以嘴裡的氣殺戮惡人。 5 公義必當他的腰帶，信實必當他脅下的帶子。

\subsubsection{耶利米書 19:11}
對他們說：『萬軍之耶和華如此說：我要照樣打碎這民和這城，正如人打碎窯匠的瓦器，以致不能再囫圇。並且人要在陀斐特葬埋屍首，甚至無處可葬。

\subsubsection{但以理書 2:44}
當那列王在位的時候，天上的神必另立一國，永不敗壞，也不歸別國的人，卻要打碎滅絕那一切國，這國必存到永遠。

\subsubsection{使徒行傳 15:28}
因為聖靈和我們定意不將別的重擔放在你們身上，唯有幾件事是不可少的，

\subsubsection{啟示錄 2:24-29，3:21，12:5，19:15，20:4 }
\begin{itemize} 
\itemsep-.25em
\item 2:24 『至於你們推雅推喇其餘的人，就是一切不從那教訓、不曉得他們素常所說撒旦深奧之理的人，我告訴你們：我不將別的擔子放在你們身上； 25 但你們已經有的，總要持守，直等到我來。 26 那得勝又遵守我命令到底的，我要賜給他權柄制伏列國， 27 他必用鐵杖轄管 b他們，將他們如同窯戶的瓦器打得粉碎，像我從我父領受的權柄一樣。 28 我又要把晨星賜給他。 29 聖靈向眾教會所說的話，凡有耳的，就應當聽！』
\item 得勝的，我要賜他在我寶座上與我同坐，就如我得了勝，在我父的寶座上與他同坐一般。
\item 3:21 婦人生了一個男孩子，是將來要用鐵杖轄管萬國的；她的孩子被提到神寶座那裡去了。
\item 19:15 有利劍從他口中出來，可以擊殺列國。他必用鐵杖轄管他們，並要踹全能神烈怒的酒榨。
\item 20:4 我又看見幾個寶座，也有坐在上面的，並有審判的權柄賜給他們。我又看見那些因為給耶穌作見證並為神之道被斬者的靈魂，和那沒有拜過獸與獸像，也沒有在額上和手上受過牠印記之人的靈魂；他們都復活了，與基督一同做王一千年。
\end{itemize} 









\subsection{陀斐特}
Tophet(地名)

字義：    焚燒之處 Place of burning，唾液 Spittle，鼓、喧聲 Drum, noise，壇 Altar，輕蔑 on tempt

彙編：    ‧王下23:10；耶7:31；賽30:33是欣嫩子谷中之高阜，以賽亞及耶利米時居民曾使子女經火奉摩洛，約西亞王禁之而污穢其地。另同名相似地乃為亞述王而設。

―― 陳瑞庭

 

陀斐特(Topheth)耶路撒冷城外、欣嫩子谷內的一個地方，以色列人曾在那裏以活人為祭獻給摩洛，褻瀆耶和華（參代下二十八3，三十三6）。

約西亞在其宗教改革中，污穢了陀斐特，並拆毀其中的祭壇（王下二十三10）。約西亞的改革似乎只起了暫時的作用，因後來耶利米對這惡行加以譴責（耶七31、32）。耶利米預言這地方要易名為「殺戮谷」，因為巴比倫人會在圍攻耶路撒冷期間，於這裏大敗猶大國（耶七32）。耶利米講述窯匠之瓦瓶的比喻時，再次重複這個預言。他強調耶路撒冷將要全然被毀，像陀斐特一樣（耶十九12）。那時，陀斐特顯然已成了一個垃圾場，打碎的瓦器會被扔在那裏，無葬身之地的屍首也會搬到陀斐特來（耶十九11）。雖然新約並沒有提及陀斐特，但「欣嫩子谷」可能就是指這個地方，意指一個遭毀壞的地方，新約中常譯作「地獄」（太五22、29、30，十28，十八9；可九43-47；路十二5）。

在加略人猶大賣主的事件中，陀斐特也可能隱晦地被提及。祭司用了猶大賣主得到的銀錢，購買窯戶的一塊田，用以埋葬外鄉人（太二十七9、10；參耶十九1-13，三十二6-9）。一個值得注意的地方是，耶利米書中稱陀斐特為「殺戮谷」，窯戶的田則稱為「血田」（太二十七10；徒一19）。然而，血田確實的位置卻不得而知。―― 證主聖經百科全書

\subsubsection{火與許多木柴早已在陀斐特為王預備 賽30:31-33}
\textcolor{blue}{亞述人必因耶和華的聲音驚惶，耶和華必用杖擊打他。 32 耶和華必將命定的杖加在他身上，每打一下，人必擊鼓彈琴。打仗的時候，耶和華必掄起手來與他交戰。} 33 原來陀斐特又深又寬，早已為王預備好了，其中堆的是火與許多木柴，耶和華的氣如一股硫磺火，使它著起來。














\subsection{窪地  בַּמְּצֻלָ֑ה  ravine,in the bottom}











\subsection{錫安山}
\subsubsection{詩篇 2:6}
說：「我已經立我的君在錫安我的聖山上了。」
\subsubsection{希伯來書 12:18-29}
18 你們原不是來到那能摸的山，此山有火焰、密雲、黑暗、暴風、 19 角聲與說話的聲音。那些聽見這聲音的，都求不要再向他們說話， 20 因為他們當不起所命他們的話說：「靠近這山的，即便是走獸，也要用石頭打死。」 21 所見的極其可怕，甚至摩西說：「我甚是恐懼戰兢。」 22 你們乃是來到錫安山，永生神的城邑，就是天上的耶路撒冷；那裡有千萬的天使， 23 有名錄在天上諸長子之會所共聚的總會，有審判眾人的神和被成全之義人的靈魂， 24 並新約的中保耶穌以及所灑的血。這血所說的比亞伯的血所說的更美。 25 你們總要謹慎，不可棄絕那向你們說話的。因為那些棄絕在地上警戒他們的，尚且不能逃罪，何況我們違背那從天上警戒我們的呢？ 26 當時他的聲音震動了地，但如今他應許說：「再一次我不單要震動地，還要震動天。」 27 這再一次的話，是指明被震動的，就是受造之物都要挪去，使那不被震動的常存。 28 所以，我們既得了不能震動的國，就當感恩，照神所喜悅的，用虔誠、敬畏的心侍奉神， 29 因為我們的神乃是烈火。
\subsubsection{啟示錄 14:1-5 }
我又觀看，見羔羊站在錫安山，同他又有十四萬四千人，都有他的名和他父的名寫在額上。 2 我聽見從天上有聲音，像眾水的聲音和大雷的聲音，並且我所聽見的好像彈琴的所彈的琴聲。 3 他們在寶座前，並在四活物和眾長老前唱歌，彷彿是新歌，除了從地上買來的那十四萬四千人以外，沒有人能學這歌。 4 這些人未曾沾染婦女，他們原是童身。羔羊無論往哪裡去，他們都跟隨他。他們是從人間買來的，做初熟的果子歸於神和羔羊。 5 在他們口中察不出謊言來，他們是沒有瑕疵的。





\subsection{小书卷}

\subsubsection{以賽亞書 8:16，28:7，29:9-12，18}
\begin{itemize}
\itemsep-.25em
\item 8:16 你要捲起律法書，在我門徒中間封住訓誨。
\item 28:7 就是這地的人，也因酒搖搖晃晃，因濃酒東倒西歪。祭司和先知因濃酒搖搖晃晃，被酒所困，因濃酒東倒西歪。他們錯解默示，謬行審判。
\item 29:9 你們等候驚奇吧！你們宴樂昏迷吧！他們醉了，卻非因酒；他們東倒西歪，卻非因濃酒。 10 因為耶和華將沉睡的靈澆灌你們，封閉你們的眼，蒙蓋你們的頭。你們的眼就是先知，你們的頭就是先見。 11 所有的默示，你們看如封住的書卷。人將這書卷交給識字的，說：「請念吧。」他說：「我不能念，因為是封住了。」 12 又將這書卷交給不識字的人，說：「請念吧。」他說：「我不識字。」
\item 29:18 那時，聾子必聽見這書上的話，瞎子的眼必從迷蒙黑暗中得以看見。
\end{itemize}

\subsubsection{耶利米書 15:16/36:1-8}
\begin{itemize}
\itemsep-.25em
\item 耶和華萬軍之神啊，我得著你的言語就當食物吃了，你的言語是我心中的歡喜快樂，因我是稱為你名下的人。
\item 猶大王約西亞的兒子約雅敬第四年，耶和華的話臨到耶利米說： 2「你取一書卷，將我對你說攻擊以色列和猶大並各國的一切話，從我對你說話的那日，就是從約西亞的日子起，直到今日，都寫在其上。 3或者猶大家聽見我想要降於他們的一切災禍，各人就回頭離開惡道，我好赦免他們的罪孽和罪惡。所以耶利米召了尼利亞的兒子巴錄來，巴錄就從耶利米口中，將耶和華對耶利米所說的一切話寫在書卷上。 5 耶利米吩咐巴錄說：「我被拘管，不能進耶和華的殿， 6 所以你要去趁禁食的日子，在耶和華殿中將耶和華的話，就是你從我口中所寫在書卷上的話，念給百姓和一切從猶大城邑出來的人聽。 7 或者他們在耶和華面前懇求，各人回頭離開惡道，因為耶和華向這百姓所說要發的怒氣和憤怒是大的。」 8 尼利亞的兒子巴錄就照先知耶利米一切所吩咐的去行，在耶和華的殿中從書上念耶和華的話。
\end{itemize}

\subsubsection{以西結書 2:8-10，3:1-11}
\begin{itemize}
\itemsep-.25em
\item 「人子啊，要聽我對你所說的話，不要悖逆，像那悖逆之家。你要開口吃我所賜給你的。」 9 我觀看，見有一隻手向我伸出來，手中有一書卷。 10 他將書卷在我面前展開，內外都寫著字，其上所寫的有哀號、嘆息、悲痛的話。
\item 1 他對我說：「人子啊，要吃你所得的，要吃這書卷，好去對以色列家講說。」 2 於是我開口，他就使我吃這書卷。 3 又對我說：「人子啊，要吃我所賜給你的這書卷，充滿你的肚腹。」我就吃了，口中覺得其甜如蜜。
他對我說：「人子啊，要吃你所得的，要吃這書卷，好去對以色列家講說。」 2 於是我開口，他就使我吃這書卷。 3 又對我說：「人子啊，要吃我所賜給你的這書卷，充滿你的肚腹。」我就吃了，口中覺得其甜如蜜。4 他對我說：「人子啊，你往以色列家那裡去，將我的話對他們講說。 5 你奉差遣不是往那說話深奧、言語難懂的民那裡去，乃是往以色列家去。 6 不是往那說話深奧、言語難懂的多國去，他們的話語是你不懂得的，我若差你往他們那裡去，他們必聽從你。 7 以色列家卻不肯聽從你，因為他們不肯聽從我，原來以色列全家是額堅心硬的人。 8 看哪，我使你的臉硬過他們的臉，使你的額硬過他們的額。 9 我使你的額像金鋼鑽，比火石更硬。他們雖是悖逆之家，你不要怕他們，也不要因他們的臉色驚惶。」 10 他又對我說：「人子啊，我對你所說的一切話，要心裡領會，耳中聽聞。 11 你往你本國被擄的子民那裡去，他們或聽或不聽，你要對他們講說，告訴他們：『這是主耶和華說的。』」
\end{itemize}

\subsubsection{但以理書 12:4-13}
但以理啊，你要隱藏這話，封閉這書，直到末時。必有多人來往奔跑，知識就必增長。」5 我但以理觀看，見另有兩個人站立，一個在河這邊，一個在河那邊。 6 有一個問那站在河水以上穿細麻衣的說：「這奇異的事到幾時才應驗呢？」 7 我聽見那站在河水以上穿細麻衣的，向天舉起左右手，指著活到永遠的主起誓說：「要到一載，二載，半載，打破聖民權力的時候，這一切事就都應驗了。」 8 我聽見這話，卻不明白，就說：「我主啊，這些事的結局是怎樣呢？」 9 他說：「但以理啊，你只管去，因為這話已經隱藏封閉，直到末時。 10 必有許多人使自己清淨潔白，且被熬煉，但惡人仍必行惡。一切惡人都不明白，唯獨智慧人能明白。 11 從除掉常獻的燔祭，並設立那行毀壞可憎之物的時候，必有一千二百九十日。 12 等到一千三百三十五日的，那人便為有福！ 13 你且去等候結局，因為你必安歇。到了末期，你必起來，享受你的福分。」

\subsubsection{撒迦利亞書 5:1-4}
我又舉目觀看，見有一飛行的書卷。 2 他問我說：「你看見什麼？」我回答說：「我看見一飛行的書卷，長二十肘，寬十肘。」 3 他對我說：「這是發出行在遍地上的咒詛，凡偷竊的必按卷上這面的話除滅，凡起假誓的必按卷上那面的話除滅。 4 萬軍之耶和華說：我必使這書卷出去，進入偷竊人的家和指我名起假誓人的家，必常在他家裡，連房屋帶木石都毀滅了。」

\subsubsection{啟 5:1-10，11:1-11}
\begin{itemize} 
\itemsep-.25em
\item 5:1 我看見坐寶座的右手中有書卷，裡外都寫著字，用七印封嚴了。 2 我又看見一位大力的天使大聲宣傳說：「有誰配展開那書卷，揭開那七印呢？」 3 在天上、地上、地底下，沒有能展開、能觀看那書卷的。 4 因為沒有配展開、配觀看那書卷的，我就大哭。 5 長老中有一位對我說：「不要哭。看哪，猶大支派中的獅子，大衛的根，他已得勝，能以展開那書卷、揭開那七印！」6 我又看見寶座與四活物並長老之中有羔羊站立，像是被殺過的，有七角七眼，就是神的七靈，奉差遣往普天下去的。 7 這羔羊前來，從坐寶座的右手裡拿了書卷。 8 他既拿了書卷，四活物和二十四位長老就俯伏在羔羊面前，各拿著琴和盛滿了香的金爐。這香就是眾聖徒的祈禱。 9 他們唱新歌，說：「你配拿書卷，配揭開七印，因為你曾被殺，用自己的血從各族、各方、各民、各國中買了人來，叫他們歸於神， 10 又叫他們成為國民做祭司，歸於神，在地上執掌王權。」
\item 我又看見另有一位大力的天使從天降下，披著雲彩，頭上有虹，臉面像日頭，兩腳像火柱。 2 他手裡拿著小書卷，是展開的。他右腳踏海，左腳踏地， 3 大聲呼喊，好像獅子吼叫。呼喊完了，就有七雷發聲。 4 七雷發聲之後，我正要寫出來，就聽見從天上有聲音說：「七雷所說的你要封上，不可寫出來。」 5 我所看見的那踏海踏地的天使向天舉起右手來， 6 指著那創造天和天上之物、地和地上之物、海和海中之物，直活到永永遠遠的，起誓說：「不再有時日了！ 7 但在第七位天使吹號發聲的時候，神的奧祕就成全了，正如神所傳給他僕人眾先知的佳音。」8 我先前從天上所聽見的那聲音又吩咐我說：「你去，把那踏海踏地之天使手中展開的小書卷取過來。」 9 我就去到天使那裡，對他說：「請你把小書卷給我。」他對我說：「你拿著吃盡了，便叫你肚子發苦，然而在你口中要甜如蜜。」 10 我從天使手中把小書卷接過來，吃盡了，在我口中果然甜如蜜；吃了以後，肚子覺得發苦了。 11 天使對我說：「你必指著多民、多國、多方、多王再說預言。」
\end{itemize} 






\subsection{鑰匙}

\subsubsection{大衛家的鑰匙-以賽亞書 22:22; 啟示錄 3:7}
\begin{itemize} 
\itemsep-.25em
\item 我必將\underline{大衛家的鑰匙}放在他肩頭上，他開無人能關，他關無人能開。
\item 「你要寫信給非拉鐵非教會的使者說：『那聖潔、真實，拿著\underline{大衛的鑰匙}，開了就沒有人能關、關了就沒有人能開的說：我知道你的行為，你略有一點力量，也曾遵守我的道，沒有棄絕我的名，看哪，我在你面前給你一個敞開的門，是無人能關的。
\end{itemize} 

\subsubsection{天國的鑰匙-馬太福音 16:19}
我要把\underline{天國的鑰匙}給你，凡你在地上所捆綁的，在天上也要捆綁；凡你在地上所釋放的，在天上也要釋放。」

\subsubsection{知識的鑰匙-路加福音 11:52}
你們律法師有禍了！因為你們把\underline{知識的鑰匙}奪了去，自己不進去，正要進去的人你們也阻擋他們。


\subsubsection{死亡和陰間的鑰匙(無底坑)-啟示錄 1:18,9:1,20:1}
\begin{itemize} 
\itemsep-.25em
\item 又是那存活的；我曾死過，現在又活了，直活到永永遠遠，並且拿著\underline{死亡和陰間的鑰匙} 
\item 第五位天使吹號，我就看見一個星從天落到地上，有\underline{無底坑的鑰匙}賜給它。它開了無底坑，便有煙從坑裡往上冒，好像大火爐的煙，日頭和天空都因這煙昏暗了
\item 我又看見一位天使從天降下，手裡拿著\underline{無底坑的鑰匙}和一條大鏈子。 2他捉住那龍，就是古蛇，又叫魔鬼，也叫撒旦，把牠捆綁一千年
\end{itemize} 




\subsection{印}
\subsubsection{神的印}
\subsubsection{兽的印记}












\subsection{預備}
\subsubsection{預備总结}
\begin{table}[H]
\centering
\begin{tabular}{l|l|l|l}\hline
\multicolumn{2}{c|}{神的預備} & \multicolumn{2}{c}{人預備的目的是迎接人子,進羔羊婚娶} \\\hline\hline
王為他兒子預備娶親的筵席& 太 22:2-14&預備人子來了& 太 24:32-44\\\hline
王為他兒子預備娶親的筵席& 路 14:15-24&預備油在器皿裡迎接新郎來進去坐席& 太 25:1-13\\\hline
預備地方& 約 14:1-3&預備永不壞的錢囊、用不盡的財寶& 路 12:33-34\\\hline
預備一座城& 來 11:16&腰束帶，燈要點,像僕人等主人從婚筵回來& 路 12:35-40\\\hline
末世要顯現的救恩& 彼前1:5&按时分粮,顺主人的意思行& 路 12:42-48\\\hline
&&預備心中盼望的緣由& 彼前 3:13-16\\\hline
&&新婦自己預備羔羊婚娶& 啟 19:7\\\hline\hline
\end{tabular}
\end{table}
 

\subsubsection{預備:迎接人子來/油/永不壞的錢囊,用不盡的財寶/腰束帶燈要點/進羔羊婚娶}
\begin{itemize}
\item \textbf{預備人子來了} - 馬太福音 24:32-44\\
你們可以從無花果樹學個比方。當樹枝發嫩長葉的時候，你們就知道夏天近了。 33 這樣，你們看見這一切的事，也該知道人子近了，正在門口了。 34 我實在告訴你們：這世代還沒有過去，這些事都要成就。 35 天地要廢去，我的話卻不能廢去。 36 但那日子、那時辰，沒有人知道，連天上的使者也不知道，子也不知道，唯獨父知道。 37 挪亞的日子怎樣，人子降臨也要怎樣。 38 當洪水以前的日子，人照常吃喝嫁娶，直到挪亞進方舟的那日， 39 不知不覺洪水來了，把他們全都沖去。人子降臨也要這樣。 40 那時，兩個人在田裡，取去一個，撇下一個； 41 兩個女人推磨，取去一個，撇下一個。 42 所以，你們要警醒，因為不知道你們的主是哪一天來到。 43 家主若知道幾更天有賊來，就必警醒，不容人挖透房屋，這是你們所知道的。 44 所以，你們也要\textcolor{red}{預備}，因為你們想不到的時候，人子就來了。
\item \textbf{預備油在器皿裡迎接新郎來進去坐席} - 馬太福音 25:1-13\\
那時，天國好比十個童女拿著燈出去迎接新郎。 2 其中有五個是愚拙的，五個是聰明的。 3 愚拙的拿著燈，卻\textcolor{red}{預備}油； 4 聰明的拿著燈，又\textcolor{red}{預備}油在器皿裡。 5 新郎遲延的時候，她們都打盹睡著了。 6 半夜有人喊著說：『新郎來了，你們出來迎接他！』 7 那些童女就都起來收拾燈。 8 愚拙的對聰明的說：『請分點油給我們，因為我們的燈要滅了。』 9 聰明的回答說：『恐怕不夠你我用的，不如你們自己到賣油的那裡去買吧。』 10 她們去買的時候，新郎到了，那\textcolor{red}{預備}好了的同他進去坐席，門就關了。 11 其餘的童女隨後也來了，說：『主啊，主啊，給我們開門！』 12 他卻回答說：『我實在告訴你們：我不認識你們。』 13 所以，你們要警醒，因為那日子、那時辰，你們不知道。
\item \textbf{預備永不壞的錢囊、用不盡的財寶} - 路加福音 12:33-34\\
你們要變賣所有的賙濟人，為自己\textcolor{red}{預備}永不壞的錢囊、用不盡的財寶在天上，就是賊不能近、蟲不能蛀的地方。 34 因為你們的財寶在哪裡，你們的心也在哪裡。

\item \textbf{腰裡束上帶，燈也要點著,好像僕人等候主人從婚姻的筵席上回來} - 路加福音 12:35-40\\
35 「你們腰裡要束上帶，燈也要點著， 36 自己好像僕人等候主人從婚姻的筵席上回來。他來到叩門，就立刻給他開門。 37 主人來了，看見僕人警醒，那僕人就有福了。我實在告訴你們：主人必叫他們坐席，自己束上帶進前伺候他們。 38 或是二更天來，或是三更天來，看見僕人這樣，那僕人就有福了。 39 家主若知道賊什麼時候來，就必警醒，不容賊挖透房屋，這是你們所知道的。 40 你們也要\textcolor{red}{預備}，因為你們想不到的時候，人子就來了。

\item \textbf{按时分粮顺主人的意思行} - 路加福音 12:42-48\\
谁是那忠心有见识的管家，主人派他管理家里的人，按时分粮给他们呢？ 43 主人来到，看见仆人这样行，那仆人就有福了。 44 我实在告诉你们：主人要派他管理一切所有的。 45 那仆人若心里说‘我的主人必来得迟’，就动手打仆人和使女，并且吃喝醉酒， 46 在他想不到的日子，不知道的时辰，那仆人的主人要来，重重地处治他，定他和不忠心的人同罪。 47 仆人知道主人的意思，却不\textcolor{red}{預備}，又不顺他的意思行，那仆人必多受责打。 48 唯有那不知道的，做了当受责打的事，必少受责打。因为多给谁，就向谁多取；多托谁，就向谁多要。

\item \textbf{預備心中盼望的緣由} - 彼得前書 3:13-16\\
你們若是熱心行善，有誰害你們呢？ 14 你們就是為義受苦，也是有福的。不要怕人的威嚇，也不要驚慌， 15 只要心裡尊主基督為聖。有人問你們心中盼望的緣由，就要常做\textcolor{red}{準備}，以溫柔、敬畏的心回答各人， 16 存著無虧的良心，叫你們在何事上被毀謗，就在何事上可以叫那誣賴你們在基督裡有好品行的人自覺羞愧。

\item \textbf{新婦自己預備羔羊婚娶} - 啟示錄 19:7\\
我們要歡喜快樂，將榮耀歸給他，因為羔羊婚娶的時候到了，新婦也自己\textcolor{red}{預備}好了，
\end{itemize}

\subsubsection{神的預備:末世要顯現的救恩/娶親的筵席}

\begin{itemize}
\item 王為他兒子預備娶親的筵席 - 馬太福音 22:2-14\\
天國好比一個王為他兒子擺設娶親的筵席， 3 就打發僕人去，請那些被召的人來赴席，他們卻不肯來。 4 王又打發別的僕人，說：『你們告訴那被召的人：我的筵席已經\textcolor{red}{預備}好了，牛和肥畜已經宰了，各樣都\textcolor{red}{齊備}，請你們來赴席。』 5 那些人不理就走了：一個到自己田裡去，一個做買賣去， 6 其餘的拿住僕人，凌辱他們，把他們殺了。 7 王就大怒，發兵除滅那些凶手，燒毀他們的城。 8 於是對僕人說：『喜筵已經\textcolor{red}{齊備}，只是所召的人不配。 9 所以你們要往岔路口上去，凡遇見的，都召來赴席。』 10 那些僕人就出去，到大路上，凡遇見的，不論善惡都召聚了來，筵席上就坐滿了客。 11 王進來觀看賓客，見那裡有一個沒有穿禮服的， 12 就對他說：『朋友，你到這裡來怎麼不穿禮服呢？』那人無言可答。 13 於是王對使喚的人說：『捆起他的手腳來，把他丟在外邊的黑暗裡，在那裡必要哀哭切齒了。』 14 因為被召的人多，選上的人少。
\item 王為他兒子預備娶親的筵席 - 路加福音 14:15-24\\
同席的有一人聽見這話，就對耶穌說：「在神國裡吃飯的有福了！」 16 耶穌對他說：「有一人擺設大筵席，請了許多客。 17 到了坐席的時候，打發僕人去對所請的人說：『請來吧，樣樣都\textcolor{red}{齊備}了。』 18 眾人一口同音地推辭。頭一個說：『我買了一塊地，必須去看看。請你准我辭了。』 19 又有一個說：『我買了五對牛，要去試一試。請你准我辭了。』 20 又有一個說：『我才娶了妻，所以不能去。』 21 那僕人回來，把這事都告訴了主人。家主就動怒，對僕人說：『快出去到城裡大街小巷，領那貧窮的、殘廢的、瞎眼的、瘸腿的來。』 22 僕人說：『主啊，你所吩咐的已經辦了，還有空座。』 23 主人對僕人說：『你出去到路上和籬笆那裡，勉強人進來，坐滿我的屋子。 24 我告訴你們：先前所請的人，沒有一個得嘗我的筵席！』」
\item 預備地方 - 約翰福音 14:1-3\\
你們心裡不要憂愁，你們信神，也當信我。 2 在我父的家裡有許多住處，若是沒有，我就早已告訴你們了，我去原是為你們\textcolor{red}{預備}地方去。 3 我若去為你們\textcolor{red}{預備}了地方，就必再來接你們到我那裡去，我在哪裡，叫你們也在哪裡。

\item 預備了一座城 - 希伯來書 11:16\\
他們卻羨慕一個更美的家鄉，就是在天上的。所以神被稱為他們的神，並不以為恥，因為他已經給他們\textcolor{red}{預備}了一座城。

\item 末世要顯現的救恩- 彼得前书1:5\\
你們這因信蒙神能力保守的人，必能得著所\textcolor{red}{預備}、到末世要顯現的救恩。

\end{itemize}










\subsection{灾难-火焰與霹雷、暴風、冰雹}
\subsubsection{以賽亞書-雷轟、地震、大聲、旋風、暴風,火焰 29:5-8}
你仇敵的群眾卻要像細塵，強暴人的群眾也要像飛糠。這事必頃刻之間，忽然臨到。 6 萬軍之耶和華必用\textcolor{red}{雷轟、地震、大聲、旋風、暴風}並吞滅的\textcolor{red}{火焰}，向她討罪。 7 那時，攻擊亞利伊勒列國的群眾，就是一切攻擊亞利伊勒和她的保障並使她困難的，必如夢景，如夜間的異象。 8 又必像飢餓的人，夢中吃飯，醒了仍覺腹空；或像口渴的人，夢中喝水，醒了仍覺發昏，心裡想喝。攻擊錫安山列國的群眾，也必如此。

\subsubsection{以賽亞書-火焰與霹雷、暴風、冰雹 30:27-30}
看哪，耶和華的名從遠方來，怒氣燒起，密煙上騰。他的嘴唇滿有憤恨，他的舌頭像吞滅的火， 28 他的氣如漲溢的河水，直漲到頸項，要用毀滅的篩籮篩淨列國，並且在眾民的口中必有使人錯行的嚼環。 29 你們必唱歌，像守聖節的夜間一樣；並且心中喜樂，像人吹笛，上耶和華的山，到以色列的磐石那裡。 30 耶和華必使人聽他威嚴的聲音，又顯他降罰的膀臂和他怒中的憤恨，並吞滅的\textcolor{red}{火焰與霹雷、暴風、冰雹}。

\subsubsection{啟示錄-火,水變為血,各樣的災殃 11:1-7}
有一根葦子賜給我當做量度的杖，且有話說：「起來，將神的殿和祭壇並在殿中禮拜的人都量一量。 2 只是殿外的院子要留下不用量，因為這是給了外邦人的，他們要踐踏聖城四十二個月。 3 我要使我那兩個見證人穿著毛衣，傳道一千二百六十天。」 4 他們就是那兩棵橄欖樹，兩個燈臺，立在世界之主面前的。 5 若有人想要害他們，就有火從他們口中出來，燒滅仇敵。凡想要害他們的都必這樣被殺。 6 這二人有權柄在他們傳道的日子叫天閉塞不下雨，又有權柄叫水變為血，並且能隨時隨意用各樣的災殃攻擊世界。 7 他們作完見證的時候，那從無底坑裡上來的獸必與他們交戰，並且得勝，把他們殺了。


\subsubsection{啟示錄-地大震動 11:11-13}
過了這三天半，有生氣從神那裡進入他們裡面，他們就站起來，看見他們的人甚是害怕。 12 兩位先知聽見有大聲音從天上來，對他們說：「上到這裡來！」他們就駕著雲上了天，他們的仇敵也看見了。 13 正在那時候，地大震動，城就倒塌了十分之一，因地震而死的有七千人，其餘的都恐懼，歸榮耀給天上的神。

\subsubsection{啟示錄-閃電、聲音、雷轟、地震、大雹 11:15-19}
第七位天使吹號，天上就有大聲音說：「世上的國成了我主和主基督的國，他要做王，直到永永遠遠。」 16 在神面前、坐在自己位上的二十四位長老就面伏於地，敬拜神， 17 說：「昔在、今在的主神，全能者啊！我們感謝你，因你執掌大權做王了。 18 外邦發怒，你的憤怒也臨到了！審判死人的時候也到了！你的僕人眾先知和眾聖徒，凡敬畏你名的人，連大帶小得賞賜的時候也到了！你敗壞那些敗壞世界之人的時候也就到了！」當時，神天上的殿開了，在他殿中現出他的約櫃，隨後有\textcolor{red}{閃電、聲音、雷轟、地震、大雹}。


\subsubsection{啟示錄-閃電、聲音、雷轟、大地震，90斤大雹子 16:12-21}
第六位天使把碗倒在幼發拉底大河上，河水就乾了，要給那從日出之地所來的眾王預備道路。 13 我又看見三個汙穢的靈，好像青蛙，從龍口、獸口並假先知的口中出來。 14 他們本是鬼魔的靈，施行奇事，出去到普天下眾王那裡，叫他們在神全能者的大日聚集爭戰。 15 （「看哪，我來像賊一樣。那警醒、看守衣服，免得赤身而行，叫人見他羞恥的有福了！」） 16 那三個鬼魔便叫眾王聚集在一處，希伯來話叫做哈米吉多頓。

第七位天使把碗倒在空中，就有大聲音從殿中的寶座上出來，說：「成了！」 18 又有\textcolor{red}{閃電、聲音、雷轟、大地震}，自從地上有人以來，沒有這樣大、這樣厲害的地震。 19 那\textcolor{red}{大城裂為三段，列國的城也都倒塌}了。神也想起巴比倫大城來，要把那盛自己烈怒的酒杯遞給他。 20 \textcolor{red}{各海島都逃避了，眾山也不見了}。 21 又有\textcolor{red}{大雹子}從天落在人身上，每一個約重一他連得(一他連得約有九十斤)。為這雹子的災極大，人就褻瀆神。

\subsubsection{啟示錄-火燒淫婦，管轄地上眾王的大城 17:6/14-18}
6我又看見那女人喝醉了聖徒的血和為耶穌作見證之人的血。我看見她，就大大地稀奇。
他們與羔羊爭戰，羔羊必勝過他們，因為羔羊是萬主之主、萬王之王。同著羔羊的，就是蒙召、被選、有忠心的，也必得勝。」 15 天使又對我說：「你所看見那淫婦坐的眾水，就是多民、多人、多國、多方。 16 你所看見的那十角與獸必恨這淫婦，使她冷落赤身，又要吃她的肉，用火將她燒盡。 17 因為神使諸王同心合意，遵行他的旨意，把自己的國給那獸，直等到神的話都應驗了。 18 你所看見的那女人就是管轄地上眾王的大城。

\subsubsection{啟示錄-巴比倫大城傾倒,死亡、悲哀、饑荒,被火燒盡 18:1-24}
此後，我看見另有一位有大權柄的天使從天降下，地就因他的榮耀發光。 2 他大聲喊著說：「巴比倫大城傾倒了！傾倒了！成了鬼魔的住處和各樣汙穢之靈的巢穴(牢獄)，並各樣汙穢可憎之雀鳥的巢穴。 3 因為列國都被她邪淫大怒的酒傾倒了，地上的君王與她行淫，地上的客商因她奢華太過就發了財。」

4 我又聽見從天上有聲音說：「我的民哪，你們要從那城出來，免得與她一同有罪，受她所受的災殃。 5 因她的罪惡滔天，她的不義，神已經想起來了。 6 她怎樣待人，也要怎樣待她，按她所行的加倍地報應她，用她調酒的杯加倍地調給她喝。 7 她怎樣榮耀自己，怎樣奢華，也當叫她照樣痛苦悲哀；因她心裡說：『我坐了皇后的位，並不是寡婦，決不至於悲哀。』 8 所以在一天之內，她的災殃要一齊來到，就是死亡、悲哀、饑荒。她又要被火燒盡了，因為審判她的主神大有能力。 9 地上的君王，素來與她行淫、一同奢華的，看見燒她的煙，就必為她哭泣哀號。 10 因怕她的痛苦，就遠遠地站著說：『哀哉！哀哉！巴比倫大城，堅固的城啊，一時之間你的刑罰就來到了！』 11 地上的客商也都為她哭泣悲哀，因為沒有人再買他們的貨物了。 12 這貨物就是金、銀、寶石、珍珠、細麻布、紫色料、綢子、朱紅色料，各樣香木，各樣象牙的器皿，各樣極寶貴的木頭和銅、鐵、漢白玉的器皿， 13 並肉桂、豆蔻、香料、香膏、乳香、酒、油、細麵、麥子、牛、羊、車、馬和奴僕、人口。

14 「巴比倫哪，你所貪愛的果子離開了你，你一切的珍饈美味和華美的物件也從你中間毀滅，決不能再見了！ 15 販賣這些貨物，藉著她發了財的客商，因怕她的痛苦，就遠遠地站著哭泣悲哀， 16 說：『哀哉，哀哉，這大城啊！素常穿著細麻、紫色、朱紅色的衣服，又用金子、寶石和珍珠為裝飾， 17 一時之間，這麼大的富厚就歸於無有了！』凡船主和坐船往各處去的，並眾水手，連所有靠海為業的，都遠遠地站著， 18 看見燒她的煙就喊著說：『有何城能比這大城呢？』 19 他們又把塵土撒在頭上，哭泣悲哀，喊著說：『哀哉，哀哉，這大城啊！凡有船在海中的都因她的珍寶成了富足，她在一時之間就成了荒場！』 20 天哪，眾聖徒、眾使徒、眾先知啊！你們都要因她歡喜，因為神已經在她身上申了你們的冤。」

21 有一位大力的天使舉起一塊石頭，好像大磨石，扔在海裡說：「巴比倫大城也必這樣猛力地被扔下去，決不能再見了。 22 彈琴、作樂、吹笛、吹號的聲音在你中間決不能再聽見。各行手藝人在你中間決不能再遇見。推磨的聲音在你中間決不能再聽見。 23 燈光在你中間決不能再照耀。新郎和新婦的聲音在你中間決不能再聽見。你的客商原來是地上的尊貴人，萬國也被你的邪術迷惑了。 24 先知和聖徒並地上一切被殺之人的血，都在這城裡看見了。」

\subsection{灾难-刀劍、饑荒、瘟疫}
\subsubsection{撒母耳記下 24:13}
於是迦得來見大衛，對他說：「你願意國中有七年的饑荒呢？是在你敵人面前逃跑，被追趕三個月呢？是在你國中有三日的瘟疫呢？現在你要揣摩思想，我好回覆那差我來的。」

\subsubsection{耶 14/15/16/21/24/27/29/32/34/38/42/43/44}
\begin{itemize}

\item 9:16 我要把他們散在列邦中，就是他們和他們列祖素不認識的列邦。我也要使刀劍追殺他們，直到將他們滅盡。」


\item 14:10-16
耶和華對這百姓如此說：「這百姓喜愛妄行 a，不禁止腳步，所以耶和華不悅納他們。現今要記念他們的罪孽，追討他們的罪惡。」 11 耶和華又對我說：「不要為這百姓祈禱求好處。 12 他們禁食的時候，我不聽他們的呼求；他們獻燔祭和素祭，我也不悅納。我卻要用刀劍、饑荒、瘟疫滅絕他們。」13 我就說：「唉，主耶和華啊！那些先知常對他們說：『你們必不看見刀劍，也不遭遇饑荒，耶和華要在這地方賜你們長久的平安。』」 14 耶和華對我說：「那些先知託我的名說假預言，我並沒有打發他們，沒有吩咐他們，也沒有對他們說話。他們向你們預言的，乃是虛假的異象和占卜，並虛無的事以及本心的詭詐。 15 所以耶和華如此說：論到託我名說預言的那些先知，我並沒有打發他們，他們還說『這地不能有刀劍饑荒』，其實那些先知必被刀劍、饑荒滅絕。 16 聽他們說預言的百姓，必因饑荒、刀劍拋在耶路撒冷的街道上，無人葬埋，他們連妻子帶兒女都是如此。我必將他們的惡倒在他們身上。

\item 15:2-3
他們問你說：『我們往哪裡去呢？』你便告訴他們：『耶和華如此說：定為死亡的，必致死亡；定為刀殺的，必交刀殺；定為饑荒的，必遭饑荒；定為擄掠的，必被擄掠。耶和華說：我命定四樣害他們，就是刀劍殺戮，狗類撕裂，空中的飛鳥和地上的野獸吞吃毀滅。

\item 16:4
他們必死得甚苦，無人哀哭，必不得葬埋，必在地上像糞土，必被刀劍和饑荒滅絕。他們的屍首必給空中的飛鳥和地上的野獸做食物。

\item 21:7
以後我要將猶大王西底家和他的臣僕百姓，就是在城內從瘟疫、刀劍、饑荒中剩下的人，都交在巴比倫王尼布甲尼撒的手中和他們仇敵並尋索其命的人手中。巴比倫王必用刀擊殺他們，不顧惜，不可憐，不憐憫。這是耶和華說的。』

\item 24:10
我必使刀劍、饑荒、瘟疫臨到他們，直到他們從我所賜給他們和他們列祖之地滅絕

\item 27:8
『無論哪一邦哪一國，不肯服侍這巴比倫王尼布甲尼撒，也不把頸項放在巴比倫王的軛下，我必用刀劍、饑荒、瘟疫刑罰那邦，直到我藉巴比倫王的手將他們毀滅。這是耶和華說的。

\item 29:16-18 
所以耶和華論到坐大衛寶座的王和住在這城裡的一切百姓，就是未曾與你們一同被擄的弟兄， 17萬軍之耶和華如此說：看哪，我必使刀劍、饑荒、瘟疫臨到他們，使他們像極壞的無花果，壞得不可吃。 18我必用刀劍、饑荒、瘟疫追趕他們，使他們在天下萬國拋來拋去，在我所趕他們到的各國中，令人咒詛、驚駭、嗤笑、羞辱。

\item 32:23-24
他們進入這地得了為業，卻不聽從你的話，也不遵行你的律法。你一切所吩咐他們行的，他們一無所行。因此，你使這一切的災禍臨到他們。 24 看哪，敵人已經來到，築壘要攻取這城，城也因刀劍、饑荒、瘟疫交在攻城的迦勒底人手中。你所說的話都成就了，你也看見了。

\item 32:36-44 現在論到這城，就是你們所說已經因刀劍、饑荒、瘟疫交在巴比倫王手中的，耶和華以色列的神如此說： 37 我在怒氣、憤怒和大惱恨中將以色列人趕到各國，日後我必從那裡將他們招聚出來，領他們回到此地，使他們安然居住。 38 他們要做我的子民，我要做他們的神。 39 我要使他們彼此同心同道，好叫他們永遠敬畏我，使他們和他們後世的子孫得福樂。 40 又要與他們立永遠的約，必隨著他們施恩，並不離開他們，且使他們有敬畏我的心，不離開我。 41 我必歡喜施恩於他們，要盡心、盡意，誠誠實實將他們栽於此地。 42 因為耶和華如此說：我怎樣使這一切大禍臨到這百姓，我也要照樣使我所應許他們的一切福樂都臨到他們。 43 你們說『這地是荒涼無人民、無牲畜，是交付迦勒底人手之地』，日後在這境內必有人置買田地。 44 在便雅憫地、耶路撒冷四圍的各處、猶大的城邑、山地的城邑、高原的城邑並南地的城邑，人必用銀子買田地，在契上畫押，將契封緘，請出見證人，因為我必使被擄的人歸回。這是耶和華說的。

\item 34：17-22 所以耶和華如此說：你們沒有聽從我，各人向弟兄、鄰舍宣告自由，看哪，我向你們宣告一樣自由，就是使你們自由於刀劍、饑荒、瘟疫之下，並且使你們在天下萬國中拋來拋去。這是耶和華說的。 18 猶大的首領，耶路撒冷的首領，太監、祭司和國中的眾民曾將牛犢劈開，分成兩半，從其中經過，在我面前立約。後來又違背我的約，不遵行這約上的話。我必將他們交在仇敵和尋索其命的人手中，他們的屍首必給空中的飛鳥和地上的野獸做食物。 21 並且我必將猶大王西底家和他的首領，交在他們仇敵和尋索其命的人與那暫離你們而去巴比倫王軍隊的手中。 22 耶和華說：我必吩咐他們回到這城，攻打這城，將城攻取，用火焚燒。我也要使猶大的城邑變為荒場，無人居住。

\item 38:2 耶和華如此說：住在這城裡的必遭刀劍、饑荒、瘟疫而死，但出去歸降迦勒底人的必得存活,就是以自己命為掠物的，必得存活。耶和華如此說：這城必要交在巴比倫王軍隊的手中,他必攻取這城。

\item 42:15-22 你們所剩下的猶大人哪，現在要聽耶和華的話！萬軍之耶和華以色列的神如此說：你們若定意要進入埃及，在那裡寄居， 16 你們所懼怕的刀劍在埃及地必追上你們，你們所懼怕的饑荒在埃及要緊緊地跟隨你們，你們必死在那裡。 17 凡定意要進入埃及在那裡寄居的，必遭刀劍、饑荒、瘟疫而死，無一人存留，逃脫我所降於他們的災禍。萬軍之耶和華以色列的神如此說：我怎樣將我的怒氣和憤怒傾在耶路撒冷的居民身上，你們進入埃及的時候，我也必照樣將我的憤怒傾在你們身上，以致你們令人辱罵、驚駭、咒詛、羞辱，你們不得再見這地方。 19 所剩下的猶大人哪，耶和華論到你們說：不要進入埃及去！你們要確實地知道，我今日警教你們了。 20 你們用詭詐自害，因為你們請我到耶和華你們的神那裡，說：『求你為我們禱告耶和華我們神，照耶和華我們的神一切所說的告訴我們，我們就必遵行。』 21 我今日將這話告訴你們，耶和華你們的神為你們的事差遣我到你們那裡說的，你們卻一樣沒有聽從。 22 現在你們要確實地知道：你們在所要去寄居之地必遭刀劍、饑荒、瘟疫而死。

\item 43:11
他要來攻擊埃及地，定為死亡的必致死亡，定為擄掠的必被擄掠，定為刀殺的必交刀殺。

\item 44:12
那定意進入埃及地在那裡寄居的，就是所剩下的猶大人，我必使他們盡都滅絕。必在埃及地仆倒，必因刀劍、饑荒滅絕，從最小的到至大的都必遭刀劍、饑荒而死，以致令人辱罵、驚駭、咒詛、羞辱。

\end{itemize}


耶利米哀歌 1:3
猶大因遭遇苦難，又因多服勞苦，就遷到外邦。她住在列國中，尋不著安息，追逼她的都在狹窄之地將她追上。
耶利米哀歌 4:15
人向他們喊著說：「不潔淨的，躲開！躲開！不要挨近我！」他們逃走漂流的時候，列國中有人說：「他們不可仍在這裡寄居。」



\subsubsection{結 5:12-17/12/14/33}
\begin{itemize}
\item 5:12-17 你的民三分之一必遭\textcolor{red}{瘟疫}而死，在你中間必因\textcolor{red}{饑荒}消滅；三分之一必在你四圍倒在\textcolor{red}{刀}下；我必將三分之一分散四方，並要拔\textcolor{red}{刀}追趕他們。

\item 6:11 主耶和華如此說：你當拍手頓足，說：『哀哉！』以色列家行這一切可憎的惡事，他們必倒在刀劍、饑荒、瘟疫之下。

\item 12:14-16 周圍一切幫助他的和他所有的軍隊，我必分散四方，也要拔刀追趕他們。…我卻要留下他們幾個人得免刀劍、饑荒、瘟疫，使他們在所到的各國中述說他們一切可憎的事，人就知道我是耶和華。

\item 14:13-23 耶和華的話臨到我說： 13 「人子啊，若有一國犯罪干犯我，我也向他伸手折斷他們的杖，就是斷絕他們的糧，使\textcolor{red}{饑荒}臨到那地，將人與牲畜從其中剪除。其中雖有挪亞、但以理、約伯這三人，他們只能因他們的義救自己的性命。這是主耶和華說的。我若使\textcolor{red}{惡獸}經過糟踐那地，使地荒涼，以致因這些獸，人都不得經過， 16 雖有這三人在其中——主耶和華說：我指著我的永生起誓，他們連兒帶女都不能得救，只能自己得救，那地仍然荒涼。或者我使\textcolor{red}{刀劍}臨到那地，說『刀劍哪，要經過那地』，以致我將人與牲畜從其中剪除， 18 雖有這三人在其中——主耶和華說：我指著我的永生起誓，他們連兒帶女都不能得救，只能自己得救。或者我叫\textcolor{red}{瘟疫}流行那地，使我滅命(帶血)的憤怒傾在其上，好將人與牲畜從其中剪除， 20 雖有挪亞、但以理、約伯在其中——主耶和華說：我指著我的永生起誓，他們連兒帶女都不能救，只能因他們的義救自己的性命。主耶和華如此說：我將這四樣大災，就是\textcolor{red}{刀劍、饑荒、惡獸、瘟疫}降在耶路撒冷，將人與牲畜從其中剪除，豈不更重嗎？然而其中必有剩下的人，他們連兒帶女必帶到你們這裡來，你們看見他們所行所為的，要因我降給耶路撒冷的一切災禍，便得了安慰。 23 你們看見他們所行所為的得了安慰，就知道我在耶路撒冷中所行的並非無故。這是主耶和華說的。


\item 33:1-6 耶和華的話臨到我說： 2 「人子啊，你要告訴本國的子民說：我使刀劍臨到哪一國，那一國的民從他們中間選立一人為守望的， 3 他見刀劍臨到那地，若吹角警戒眾民， 4 凡聽見角聲不受警戒的，刀劍若來除滅了他，他的罪(血)就必歸到自己的頭上。 5 他聽見角聲，不受警戒，他的罪必歸到自己的身上；他若受警戒，便是救了自己的性命。 6 倘若守望的人見刀劍臨到，不吹角，以致民不受警戒，刀劍來殺了他們中間的一個人，他雖然死在罪孽之中，我卻要向守望的人討他喪命的罪 b。

\item 33:27-29 你要對他們這樣說：『主耶和華如此說：我指著我的永生起誓，在荒場中的必倒在\textcolor{red}{刀}下，在田野間的必交給\textcolor{red}{野獸}吞吃，在保障和洞裡的必遭\textcolor{red}{瘟疫}而死。我必使這地荒涼，令人驚駭，她因勢力而有的驕傲也必止息。以色列的山都必荒涼，無人經過。 29 我因他們所行一切可憎的事，使地荒涼，令人驚駭，那時他們就知道我是耶和華。』


\end{itemize}


\subsubsection{何西阿書 9:6}
看哪，他們逃避災難，埃及人必收殮他們的屍首，摩弗人必葬埋他們的骸骨。他們用銀子做的美物上必長蒺藜，他們的帳篷中必生荊棘

\subsubsection{阿摩司書 4:6-12}
「我使你們在一切城中牙齒乾淨，在你們各處糧食缺乏，你們仍不歸向我。」這是耶和華說的。
阿摩司書 9:4
雖被仇敵擄去，我必命刀劍殺戮他們。我必向他們定住眼目，降禍不降福。」

撒迦利亞書 7:14
我必以旋風吹散他們到素不認識的萬國中。這樣，他們的地就荒涼，甚至無人來往經過，因為他們使美好之地荒涼了。」

\subsubsection{路加福音 21:24}
他們要倒在刀下，又被擄到各國去。耶路撒冷要被外邦人踐踏，直到外邦人的日期滿了。


\subsubsection{啟示錄 6:4-8}
就另有一匹馬出來，是紅的，有權柄給了那騎馬的，可以從地上奪去太平，使人彼此相殺，又有一把大刀賜給他。
我就觀看，見有一匹灰色馬。騎在馬上的名字叫做「死」，陰府也隨著他。有權柄賜給他們，可以用刀劍、饑荒、瘟疫、野獸殺害地上四分之一的人。














